% =====================================================================
% EgoAnchor — IEEE VR 2027 最终版主稿（中文）
% 历史草稿：egoanchor_cn_v1.tex / egoanchor_cn_outline.tex（保留备查，勿覆盖）
% 编译产物输出到 2026-EgoAnchor/pdf/
% 参考文献：egoanchor_cn_refs.bib
% =====================================================================
\documentclass[review]{vgtc}

\graphicspath{{figures/}{pictures/}{images/}{./}}

\usepackage[UTF8]{ctex}
\usepackage{times}
\renewcommand*\ttdefault{txtt}
\usepackage{mathptmx}
\usepackage{booktabs}
\usepackage{tabularx}
\usepackage{multirow}
\usepackage{enumitem}
\usepackage{amsmath}
\usepackage{amssymb}
\usepackage{graphicx}
\usepackage{xcolor}

\onlineid{0}
\vgtccategory{Research}
% vgtc 模板在此统一插入 hyperref / cleveref 等兼容包，
% 不要再手动 \usepackage{hyperref} 或 \usepackage{cleveref}。
\vgtcinsertpkg

\title{EgoAnchor：面向开放消费级混合现实的动态真实物体锚定}

\author{作者待定}
\authorfooter{
% 占位：匿名审稿阶段使用 anonymous；正式投稿时填入作者信息
}

\teaser{
  \centering
  \fbox{\rule{0pt}{2.0in}\rule{0.95\linewidth}{0pt}}
  \caption{EgoAnchor 系统概览。}
  \label{fig:teaser}
}

\abstract{%
动态真实物体锚定是实现以真实物体为中心的透视混合现实（Passthrough Mixed Reality，PMR）交互的基础能力。然而，现有方案通常依赖物理标签、专用深度传感器、封闭平台接口或逐物体离线准备流程，因而难以部署到开放消费级场景中的日常物体。本文提出 EgoAnchor，一套面向开放消费级混合现实的动态真实物体锚定系统。该系统仅依赖头戴显示器双目图像与目标物体的可用三维模型，无需物理标签、专用深度硬件或逐物体离线训练，即可为广泛的日常刚性物体提供连续6DoF动态锚定能力。为此，我们构建了一套由视觉感知后端与对象锚定运行时组成的分层架构：前者负责持续生成异步物体位姿流，后者将异步位姿流持续转换为世界坐标系下稳定、连续且可恢复的对象锚点。我们在 Meta Quest 3 上完成系统实现，并通过定量实验和用户研究评估系统在锚定精度、稳定性、响应行为以及真实交互中的表现。结果表明，EgoAnchor 能够在开放消费级混合现实场景中有效支持动态真实物体锚定，并在锚定稳定性与交互体验方面展现出显著优势。}

\keywords{透视混合现实, 动态对象锚定, 6DoF位姿估计, 动态物体追踪, 6D物体定位}

\begin{document}

\maketitle

\section{引言}

透视混合现实（Passthrough Mixed Reality, PMR）通过将虚拟内容实时叠加到真实环境，使用户能够直接围绕现实世界进行观察、操作与交互。近年来，随着 Meta Quest 3、Apple Vision Pro 等消费级头戴显示设备的普及，越来越多的混合现实应用开始围绕真实物体展开，例如虚拟说明书附着、物体高亮引导、桌面协同办公以及近物精细交互等。这类应用共同依赖于此项基础能力：使虚拟内容能够持续、稳定地附着于真实物体，并随物体运动保持一致的空间关系。

对于静态环境，这一目标通常可以通过空间锚点（Spatial Anchor）实现。然而，当用户拿起、移动、旋转或传递真实物体时，固定于环境中的空间锚点无法反映物体自身的运动，虚拟内容将随之产生明显偏移。因此，混合现实系统需要建立能够随真实物体运动持续更新的对象级空间参考，即动态对象锚定（Dynamic Object Anchoring）：无论用户自身运动、物体发生位姿变化，还是经历短暂遮挡，虚拟内容都应始终与真实物体保持稳定一致的六自由度（6DoF）对应关系。

近年来，消费级混合现实平台开始提供免标记对象追踪能力。HoloLens~2 的  Azure Object Anchors 利用设备感知与对象数字模型完成对象匹配，但公开资料显示其能力主要面向较大尺寸的静态对象；Apple Vision Pro 提供对象追踪能力，但依赖基于三维模型构建的 reference object 以及 Create ML 的离线训练流程；Vuforia Model Targets 基于对象形状以及 CAD/3D 模型实现真实物体识别与跟踪，但因其受约束的初始化与准备流程，常见于工业与设备场景；Meta Quest 近年来也开始支持动态对象追踪，但当前仍属于受系统支持对象集合与用户设置约束的受控能力。总体而言，现有平台主要提供封装好的对象追踪能力，通常依赖专用硬件、封闭的平台接口、模型准备或对象级离线流程，难以直接支持开放场景中的广泛日常物体，也较少向开发者开放底层视觉感知结果。

为此，本文提出 EgoAnchor，一套面向开放消费级混合现实的动态真实物体锚定系统。该系统采用对象感知与对象锚定解耦的分层架构，由视觉感知后端（Visual Perception Backend）和对象锚定运行时（Object Anchoring Runtime）组成。前者负责从头戴设备双目图像中持续生成异步物体位姿流，后者进一步将异步视觉位姿流转换为世界坐标系下稳定、连续且可恢复的对象锚点，从而为混合现实应用提供统一的动态对象锚定能力。整个系统仅依赖头戴显示器双目图像与目标物体的可用三维模型，无需物理标记、专用深度设备或逐物体离线训练，并能够与近年来的 Image-to-3D 技术结合，为开放场景中的更多日常刚性物体提供部署潜力。

我们在 Meta Quest 3 上实现了 EgoAnchor，并通过定量实验和用户研究对系统进行了评估。定量实验利用官方控制器提供的高精度位姿作为参考，测量系统在锚定精度、稳定性与响应行为方面的表现；用户研究围绕典型混合现实交互任务展开，评估系统在真实应用中的可用性。实验结果表明，EgoAnchor 能够在开放消费级混合现实场景中稳定支持动态真实物体锚定，并为围绕真实物体展开的交互提供更一致的空间附着与更好的整体体验。

本文的主要贡献如下：

\begin{enumerate}

\item 提出 EgoAnchor，一套面向开放消费级混合现实的动态真实物体锚定系统。该系统仅依赖头戴显示器双目图像与目标物体的可用三维模型，无需物理标记、专用深度硬件或逐物体离线训练，能够为广泛的日常刚性物体提供连续6DoF动态锚定能力。

\item 设计一套由视觉感知后端与对象锚定运行时组成的分层系统架构：视觉感知后端负责从头戴设备双目图像中持续生成统一的异步物体位姿流；对象锚定运行时则将异步视觉位姿流持续转换为世界坐标系下稳定、连续且可恢复的对象锚点。

\item 在 Meta Quest 3 上完成系统实现，并通过定量实验与用户研究评估系统在锚定精度、稳定性、响应行为以及真实交互体验方面的表现。

\end{enumerate}

\section{相关工作}

\subsection{平台级真实物体锚定}

真实物体锚定旨在为物理对象建立稳定的对象级空间参考，使虚拟内容能够持续附着于真实物体，是近年来混合现实平台的重要能力。随着消费级头戴显示设备的发展，主流 XR 平台逐步开始提供对象追踪或对象锚定能力，但在开放部署、对象泛化以及动态支持方面仍存在一定局限。

早期系统主要采用基于物理标记的方法，通过 AprilTag、ArUco 等人工标记获得稳定的位姿估计，ARToolKit、ARCore Augmented Images 等系统均属于这一范式。这类方法实现简单、鲁棒性较高，但需要预先在物体表面贴附标记，不仅改变真实物体外观，也难以适用于开放场景中广泛的日常刚性物体。

近年来，消费级混合现实平台开始提供免标记对象追踪能力。HoloLens~2 的 Azure Object Anchors 利用主动深度感知完成对象锚定；Apple Vision Pro 提供对象追踪能力，但仍依赖 LiDAR 扫描以及 Create ML 的离线对象训练；Vuforia Model Targets 利用 CAD 模型完成对象识别，主要面向静态工业目标；Meta Quest 近年来也开始支持动态对象追踪，但目前仍局限于少量预定义对象。

总体而言，现有平台主要提供封装好的对象追踪能力，通常依赖专用硬件、封闭的平台接口或对象级离线准备，难以支持开放场景中广泛的日常刚性物体，也较少向开发者开放底层视觉感知结果。因此，如何利用开放视觉感知能力，在消费级设备上构建面向开放场景的动态真实物体锚定系统，仍然具有重要研究价值。

\subsection{6DoF物体位姿估计与跟踪}

6DoF物体位姿估计旨在从图像中恢复刚性物体相对于相机的空间位姿，是动态真实物体锚定的感知基础。近年来，随着视觉基础模型的发展，相关研究不断降低对逐物体训练和专用硬件的依赖，使开放场景中的对象感知成为可能。

早期方法通常依赖局部几何特征与 PnP 求解，在弱纹理、遮挡以及复杂光照条件下鲁棒性有限。随后，大量深度学习方法针对特定物体学习位姿估计网络，如 PoseCNN、DenseFusion 等。进一步地，CosyPose、MegaPose 等方法结合渲染匹配提升了未知物体的泛化能力，FoundationPose 则实现了无需逐物体训练的零样本6DoF位姿估计。随着开放词表目标检测、视频对象分割、立体匹配等视觉基础模型的发展，开放场景中的对象感知能力得到进一步提升。然而，这些能力大多以独立视觉模块的形式存在，如何将其组织为面向混合现实应用的统一视觉感知流水线，仍然缺乏系统性的研究。

然而，上述工作的目标始终是提升相机坐标系下6DoF位姿估计的准确性，而非混合现实中的对象锚定。视觉推理产生的异步时延、估计噪声以及跟踪波动，使高精度位姿流仍难以直接驱动稳定的虚实配准。因此，现有研究主要关注对象感知（Object Perception），而如何将异步视觉位姿流进一步转换为世界坐标系下稳定、连续且可恢复的对象锚点，仍然缺乏系统性的研究。

\subsection{XR 时延补偿与时间一致性}

配准误差（Registration Error）与端到端时延（End-to-end Latency）自增强现实诞生以来便被认为是影响用户体验的核心问题，也是导致虚实错位、视觉疲劳与沉浸感下降的重要原因。Azuma 的经典综述指出，时延与运动速度的乘积是动态配准误差的主要来源~\cite{azuma1997survey}，Holloway 的系统分析进一步量化了不同误差源对配准精度的影响~\cite{holloway1997registration}。对于依赖异步视觉感知的混合现实系统而言，图像采集、网络传输与模型推理通常会引入数十毫秒乃至更高的端到端时延；当物体或用户快速运动时，这种异步时延便会表现为虚拟内容相对于真实物体的明显漂移与打滑。

围绕时延补偿，学界与工业界发展出两类互补思路。其一是预测式跟踪：Azuma 与 Bishop 通过外推头部位姿，将动态配准误差降低数倍~\cite{azuma1994improving}；双指数平滑等方法提供了更轻量的预测替代~\cite{laviola2003double}。在交互输入侧，One~Euro 滤波器以速度自适应的方式权衡抖动与延迟~\cite{casiez2012oneeuro}，是三维用户界面处理含噪输入的常用基线，但也揭示了卡尔曼滤波~\cite{kalman1960new} 在快速运动下引入滞后的固有折衷。其二是后渲染对齐：Mark 等人的后渲染三维 warp 将已渲染画面重投影到新视角~\cite{mark1997postrendering}，成为时间扭曲的概念源头；van~Waveren 的异步时间扭曲（ATW）在扫描输出前用最新头显位姿重投影画面~\cite{vanwaveren2016atw}，已成为消费级头显的标准时延补偿手段。

上述工作主要围绕头部运动预测、渲染重投影以及交互输入滤波展开，其目标是降低显示链路中的时延影响。相比之下，由异步视觉模块输出的对象位姿对应于历史采集时刻，其世界坐标关系需要结合设备历史位姿进行恢复。尽管 OpenXR 及各类平台运行时提供了设备级时空同步能力，这些机制主要面向设备自身位姿管理，尚未进一步支持由第三方异步视觉感知驱动的动态真实物体锚定。因此，如何面向异步视觉位姿流建立统一的对象锚定运行时，仍然是消费级混合现实系统中的关键问题。

\section{系统设计}

\subsection{设计目标与系统挑战}

EgoAnchor 的目标是将异步外部感知流转化为应用层可直接使用的真实物体锚点。系统需要同时满足四个设计目标：

\noindent\textbf{(1)~世界一致性（World Consistency）。} 锚点应对应观测时刻的正确世界位置，而非回包到达时刻的近似位置。这要求系统解决\emph{时间异步}挑战：设采集时刻为 $t_i$、外部追踪包返回时刻为 $t_r$，外部感知给出相机坐标系观测 ${}^{C}\mathbf{T}_{O}(t_i)$，头显定位给出采集时刻相机世界位姿 ${}^{W}\mathbf{T}_{C}(t_i)$。若直接使用到达时刻映射 ${}^{W}\hat{\mathbf{T}}_{O} = {}^{W}\mathbf{T}_{C}(t_r)\cdot{}^{C}\mathbf{T}_{O}(t_i)$，因 ${}^{W}\mathbf{T}_{C}(t_r)\neq{}^{W}\mathbf{T}_{C}(t_i)$ 而引入打滑，其量级随头部运动速度与时延 $(t_r-t_i)$ 的乘积增长。

\noindent\textbf{(2)~稳定可用性（Stability）。} 锚点更新应抑制抖动、跳变与明显错误更新。外部视觉追踪在运动模糊、光照变化或局部遮挡下不可避免地会输出噪声乃至位姿跃变，这要求系统具备\emph{质量感知的门控机制}。

\noindent\textbf{(3)~可恢复性（Recoverability）。} 目标被遮挡、出视野或短时跟踪失败后，系统应能恢复。这要求设计\emph{生命周期管理}以处理追踪中断与重获取。

\noindent\textbf{(4)~可分析性（Diagnosability）。} 系统应暴露足够的状态、时延与诊断量，以支持实验评估与问题排查。

\subsection{系统架构概览}

EgoAnchor 采用前后端分离的分布式架构（见图~\ref{fig:system-architecture}，占位），由四个环节构成：

\noindent\textbf{前端采集（Quest/Unity）。} Quest~3 头显以 60\,fps 采集双目透视图像，每帧绑定唯一的 \texttt{frame\_id} 与采集时刻的相机世界位姿 ${}^{W}\mathbf{T}_{C}(t_i)$，写入环形缓存。图像与相机参数通过 ZMQ Protobuf 数据面以 PUB/SUB 模式发布。

\noindent\textbf{后端追踪（Python服务）。} 运行于消费级 GPU（RTX~5080 Laptop 约 5\,fps，RTX~5090 Desktop 约 12\,fps）的追踪服务订阅图像流，依次执行：(i)~基于 prompt 的目标级分割（YOLOE-26~\cite{yoloe2025} / SAM3~\cite{sam3_2025}），(ii)~双目立体深度估计（FoundationStereo~\cite{foundationstereo2025}），(iii)~基于 FoundationPose~\cite{foundationpose2023} 的 register / track / re-register。输出相机坐标系位姿 ${}^{C}\mathbf{T}_{O}(t_i)$ 与多维可靠性评分，通过 NATS Protobuf 消息面发布。

\noindent\textbf{四层锚定运行时（Unity）。} 订阅后端位姿流，依次经过：(1)~时间对齐层（帧对齐锚定），(2)~质量门控层（可靠性评分过滤），(3)~时序平滑层（卡尔曼状态估计与静止锁定），(4)~生命周期层（保持/重获取决策）。四层输出的锚点更新以后端感知频率到达（5--12\,fps），系统通过卡尔曼预测在渲染帧间进行时间升采样，最终以 60\,fps 驱动应用层虚拟内容，确保视觉流畅性。

\noindent\textbf{命令与控制（NATS Request/Reply）。} Unity 通过 NATS 命令面向后端发送重获取请求，后端回复执行状态，形成闭环自愈。

\section{四层协同锚定架构}

EgoAnchor 的核心是将异步外部感知流转化为应用层可用锚点的四层协同架构。我们首先介绍感知前端（§4.1），随后详细阐述四层锚定运行时（§4.2--§4.5）。

\subsection{对象级感知前端}

\noindent\textbf{相机参数映射。} Quest~3 Passthrough API 提供双目图像（原始分辨率 1440$\times$1920）与相机内参。系统将图像下采样至处理分辨率（如 720$\times$960），并相应调整内参矩阵的焦距与主点。

\noindent\textbf{目标级分割。} 采用基于 prompt 的语义分割（YOLOE-26~\cite{yoloe2025} 或 SAM3~\cite{sam3_2025}），用户通过自然语言或点击指定目标物体，系统输出双目掩膜 $M_L, M_R$。

\noindent\textbf{双目深度估计。} FoundationStereo~\cite{foundationstereo2025} 在双目掩膜区域内估计度量深度图 $D_L$，为后续 6DoF 估计提供尺度信息。

\noindent\textbf{Register / Track / Re-register。} 系统维护 register（初始化）、track（在线追踪）、re-register（重获取）三种模式。Register 与 re-register 使用 FoundationPose~\cite{foundationpose2023} 的完整优化，track 使用轻量迭代更新。输出相机坐标系位姿 ${}^{C}\mathbf{T}_{O}(t_i)$ 与可靠性评分（颜色重投影误差 $\mathit{score}_{\text{reproj}}$、渲染深度有效性 $\mathit{score}_{\text{depth}}$）。

\subsection{第一层：时间对齐——帧对齐锚定}

时间对齐层解决异步感知的时间错配问题。传统到达时刻映射直接使用当前最新头显位姿，导致头动打滑。EgoAnchor 采用\textbf{帧对齐锚定（frame-aligned anchoring）}：

\noindent\textbf{机制。} Unity 前端为每帧透视图像分配唯一 \texttt{frame\_id}，并将采集时刻的相机世界位姿 ${}^{W}\mathbf{T}_{C}(t_i)$ 写入环形缓存（容量约 300 帧，覆盖 5 秒历史）。当后端返回对应 \texttt{frame\_id} 的相机坐标系位姿 ${}^{C}\mathbf{T}_{O}(t_i)$ 时，Unity 按 \texttt{frame\_id} 回查采集时刻的缓存位姿，计算世界空间映射：
\begin{equation}
  {}^{W}\mathbf{T}_{O}(t_i) = {}^{W}\mathbf{T}_{C}(t_i) \cdot {}^{C}\mathbf{T}_{O}(t_i).
  \label{eq:frame-aligned}
\end{equation}

\noindent\textbf{与时间扭曲的对偶关系。} 该机制可视为渲染管线 late latching~\cite{vanwaveren2016atw,mark1997postrendering} 在感知融合层的对偶：时间扭曲把已渲染画面对齐到\emph{当前}最新头显位姿（补偿渲染时延），帧对齐锚定把外部感知结果对齐回其\emph{采集时刻}头显位姿（补偿感知时延）。二者都通过时间标签对齐消除运动误差，而不依赖预测。

\noindent\textbf{坐标系对齐。} Quest SDK 输出左目/右目相机位姿，需选定参考相机（如左目）并补偿轴约定差异（OpenGL vs Unity）。

\subsection{第二层：质量门控——可靠性评分过滤}

质量门控层拦截异常观测，防止低质量位姿污染锚点。

\noindent\textbf{多维可靠性评分。} 后端计算两类评分：(i)~颜色重投影误差 $\mathit{score}_{\text{reproj}} \in [0,1]$，度量渲染三维模型与当前透视图像在 LAB 色彩空间下的相似度；(ii)~渲染深度有效性 $\mathit{score}_{\text{depth}} \in [0,1]$，度量掩膜内有效深度覆盖比例。几何对数平均：
\begin{equation}
  S_t = \exp\!\left(w_c\ln(\mathit{score}_{\text{reproj}}) + w_d\ln(\mathit{score}_{\text{depth}})\right),
  \label{eq:reliability}
\end{equation}
权重采用 $w_c=0.2, w_d=0.8$（深度一致性比颜色更可靠）。

\noindent\textbf{门控逻辑。} 仅当 $S_t \geq \theta_{\text{gate}}$（如 0.6）且追踪状态为 Tracking 时，观测进入下游时序平滑层；否则触发保持或重获取（§4.5）。

\subsection{第三层：时序平滑——卡尔曼状态估计与静止锁定}

时序平滑层将通过门控的观测转化为平滑、稳定的锚点输出。

\noindent\textbf{恒速卡尔曼状态估计。} 采用恒速模型~\cite{kalman1960new} 对物体世界坐标位置 $\mathbf{p}_t$、速度 $\mathbf{v}_t$ 与旋转四元数对数空间 $\mathbf{r}_t$ 进行自适应估计。状态方程：
\begin{equation}
  \mathbf{x}_{t+1} = \mathbf{F}\mathbf{x}_t + \mathbf{w}_t, \quad \mathbf{z}_t = \mathbf{H}\mathbf{x}_t + \mathbf{v}_t,
\end{equation}
其中 $\mathbf{x}_t = [\mathbf{p}_t, \mathbf{v}_t, \mathbf{r}_t]^\top$。过程噪声 $\mathbf{Q}$ 与观测噪声 $\mathbf{R}$ 根据可靠性评分自适应调整。

\noindent\textbf{时间升采样。} 后端感知以低频（5--12\,fps）到达，而头显渲染需要 60\,fps 以上的锚点更新以保持视觉流畅。系统支持两种升采样策略：(i)~延迟插值（DelayedInterp）主动牺牲 1--2 个观测周期（约 100--200\,ms）的延迟，在已到达的控制点之间进行样条插值（Hermite / Catmull-Rom），保证无 overshoot 与轨迹平滑；(ii)~误差融合（Blend）零延迟外推到当前时刻（有安全限制），新观测到达时不直接跳变，而是计算残差并在后续帧指数衰减，实现 C⁰ 连续与平滑收敛。两种策略都依赖运动模型（Kalman / One Euro / Raw）提供控制点与预测接口，在精度优先与响应优先间提供灵活选择。

\noindent\textbf{静止锁定（Static Lock）。} 当物体与头显相对静止时，即使经过卡尔曼平滑，仍会残留微小抖动。静止锁定通过以下机制完全抑制噪声：

\begin{itemize}[leftmargin=1.4em]
  \item \textbf{进入条件：} 连续观测速度和角速度低于极低门槛（如 0.01\,m/s、1\,deg/s）且持续时间达到 $\tau_{\text{settle}}$（如 0.5\,s），锁定输出为进入时刻的世界位姿 $\mathbf{T}_{\text{locked}}$。
  \item \textbf{解锁判定：} 三路证据累积触发解锁：(i)~速度逃逸（卡尔曼估计速度超出阈值），(ii)~漂移绳索（观测共识点与锁定时刻锚点的距离超出阈值），(iii)~CUSUM 偏差累积（残差累积和超出死区）。
  \item \textbf{接缝消除：} 解锁瞬间，锁定输出与当前卡尔曼估计间存在接缝残差 $\Delta\mathbf{T}_{\text{seam}}$，系统在独立的接缝消除阶段将其指数衰减至零，平滑过渡回正常滤波状态。
\end{itemize}

\subsection{第四层：生命周期——保持与重获取}

生命周期层管理追踪中断与恢复，系统维护 Searching、Tracking、Coasting、Lost 四种状态：

\noindent\textbf{Searching $\rightarrow$ Tracking。} 初始化或重获取成功后，可靠性评分连续高于阈值，进入 Tracking。

\noindent\textbf{Tracking $\rightarrow$ Coasting。} 可靠性评分低于阈值（如遮挡）时，进入 Coasting。系统保持卡尔曼预测输出，但不更新观测；虚拟内容保持可见但带有视觉提示（如半透明）。

\noindent\textbf{Coasting $\rightarrow$ Tracking / Lost。} 若在 $\tau_{\text{coast}}$（如 2\,s）内恢复高质量观测，回到 Tracking；否则进入 Lost，虚拟内容隐藏。

\noindent\textbf{Lost $\rightarrow$ Searching（重获取）。} Unity 通过 NATS 命令面向后端发送 re-register 请求，后端执行完整的 FoundationPose 初始化，成功后回到 Searching。

这四层协同设计将 5--12\,fps 的异步外部感知流转化为 60\,fps 的世界一致、稳定可恢复的锚点输出。

\section{实现}

\subsection{硬件与软件环境}
% 写作指引：Quest 3、RTX 5080 Laptop(~5fps)/5090 Desktop(~12fps)、Unity、NATS 等；强调消费级硬件与升采样到 60fps。
给出可复现的设备与软件细节，并说明从 5--12\,fps 异步感知到 60\,fps 输出的时序升采样。

\subsection{网络与通信层实现}
% 写作指引：Protobuf 序列化；ZMQ PUB/SUB 数据面 + NATS Core 消息面 + request/reply 命令面。
说明双平面/三语义通道的运行时实现。

\subsection{离线仿真与诊断支持}
% 写作指引：Tools3 离线 JSONL 回放与参数敏度分析。
介绍离线升采样仿真与诊断量，支撑实验分析与复现。

\section{评估设计}

\subsection{研究问题}

本文评估的对象是 EgoAnchor 作为\textbf{端到端系统}的整体性能，而不仅是单个模块的算法新颖性。我们聚焦以下四个研究问题：

\begin{itemize}[leftmargin=1.4em]
  \item \textbf{RQ1（系统整体性能）：} EgoAnchor 能否在消费级硬件上将异步外部 6DoF 感知流转化为世界一致、稳定可用的真实物体锚点？与平台级方案的定性对比如何？
  \item \textbf{RQ2（四层架构贡献）：} 四层协同架构（时间对齐、质量门控、时序平滑、生命周期管理）在消除打滑、抑制抖动、降低延迟上的消融收益如何？
  \item \textbf{RQ3（鲁棒性与恢复）：} 系统在遮挡、出视野、快速头动等干扰条件下的恢复成功率、恢复时间与失效模式如何？
  \item \textbf{RQ4（任务可用性）：} 稳定的真实物体锚定能否提升用户在对齐与恢复任务中的客观表现与主观信任？
\end{itemize}

\subsection{基线与消融组合}

为全面评估系统性能，我们设计两层次对比：

\noindent\textbf{系统级对齐基线（回答 RQ1）：}
\begin{itemize}[leftmargin=1.4em]
  \item \textbf{Arrival-Time Mapping：} 传统到达时刻映射，无帧对齐，用于证明时间对齐的必要性。
  \item \textbf{Frame-Aligned Raw：} 仅使用帧对齐，无门控与时序处理，展现视觉追踪的原始精度与噪声。
  \item \textbf{EgoAnchor (Default)：} 完整四层架构（帧对齐 + 门控 + 卡尔曼 + 静止锁定 + 生命周期管理），系统推荐配置。
\end{itemize}

\noindent\textbf{模块化消融（回答 RQ2）：}
\begin{itemize}[leftmargin=1.4em]
  \item \textbf{时序滤波器对比：} 在保持帧对齐与门控不变的前提下，对比 Low-Pass、One Euro~\cite{casiez2012oneeuro}、Kalman~\cite{kalman1960new} 在抖动-延迟权衡上的表现。\textbf{重要：}须扫描各滤波器参数（如 One Euro 的 \texttt{fcmin}/\texttt{beta}、Kalman 的 $\mathbf{Q}$、Low-Pass 的截止频率）画出 Pareto 前沿，并标注实际部署工作点，避免单点比较的不可辩护性。
  \item \textbf{静止锁定消融：} 对比启用/禁用静止锁定在物体与头显相对静止时的微抖动抑制效果。
  \item \textbf{质量门控消融：} 在遮挡、低光等干扰场景下，对比有/无可靠性门控在防止位姿跃变与触发重获取中的保护作用。
\end{itemize}

\subsection{测量指标}

\noindent\textbf{主指标（锚点质量）：}
\begin{itemize}[leftmargin=1.4em]
  \item \textbf{世界锚定误差：} 虚拟 Overlay 在 Unity 世界坐标中的 6DoF 位姿与 Ground Truth 之间的平移 RMSE（米）与旋转测地 RMSE（度），报告均值与 95 分位。
  \item \textbf{静态抖动：} 物体与头显均静止时，锚点输出的标准差（世界空间，毫米级）。
  \item \textbf{头动打滑：} 头显以不同速度运动时，虚拟 Overlay 相对真实物体的视觉滑动位移（毫米），绘制打滑随头速曲线。
\end{itemize}

\noindent\textbf{辅助指标（系统性能）：}
\begin{itemize}[leftmargin=1.4em]
  \item \textbf{端到端时延：} 从图像采集到锚点更新的总时延，采用 Ground Truth 轨迹与渲染轨迹的互相关时滞分析~\cite{diluca2010latency}，报告 P50/P90。
  \item \textbf{恢复时间：} 从遮挡移除到系统重新稳定锚定的时间，须报告随阈值变化的曲线而非单一数值。
  \item \textbf{恢复成功率：} 在规定时间窗口内（如 3\,s）成功重获取的比例。
\end{itemize}

\noindent\textbf{支持性指标（感知质量）：}相机坐标系位姿误差（平移/旋转、ADD/ADD-S~\cite{hinterstoisser2012add,posecnn2017}）作为支持性指标说明底层感知质量，但不作为正文评价主线。

\subsection{Ground Truth 获取与实验条件}

\noindent\textbf{真值获取。} 主定量基准采用 \textbf{Quest Touch 手柄}作为标靶：其 CAD 模型已知，6DoF 世界位姿可由 Meta SDK 原生获取。在运行中，EgoAnchor 仅使用 Passthrough 图像对手柄进行视觉追踪，而 SDK 输出作为绝对 Ground Truth，用于高精度计算世界锚定误差与头动打滑。对于日常物体（键盘、鼠标等），在物体隐蔽处粘贴 ArUco 标记~\cite{garridojurado2014aruco}或 AprilTag~\cite{olson2011apriltag}，由独立标定相机离线解算真值，EgoAnchor 运行时对标记不可见。

\noindent\textbf{须披露的方法学局限：}(i)~手柄真值与锚点共用 Quest 跟踪系，只能隔离锚定/滤波栈，无法暴露与头显 SLAM 同源的共模误差；(ii)~单标记近正对存在旋转翻转歧义（IPPE flip），须改用多标记板并报告真值静态重复性作为噪声地板，只在噪声地板之上断言误差；(iii)~与 Vision Pro 的跨设备对比仅作\emph{定性}对照，须共享同一物理标记板并显式报告对齐残差。

\noindent\textbf{实验条件。} 覆盖以下代表性场景：
\begin{enumerate}[leftmargin=1.4em]
  \item \textbf{静态观察：}物体与头显均静止，测量静态抖动与长时漂移。
  \item \textbf{快速头动：}头显以不同速度（慢速 20\,deg/s、中速 60\,deg/s、快速 120\,deg/s）转动，测量头动打滑与锚定误差。
  \item \textbf{部分遮挡：}用书本或手遮挡物体 50\%--80\% 区域，测量门控响应与恢复时间。
  \item \textbf{出视野：}物体移出视野后重新观察，测量重获取成功率与恢复时间。
\end{enumerate}

\subsection{用户研究设计}

\noindent\textbf{设计。} 被试内（within-subject）设计，N=12--15（报告先验功效分析，G*Power 计算），拉丁方平衡条件顺序。两条件：\textbf{Baseline}（Frame-Aligned + Kalman，无静止锁定与门控）vs \textbf{EgoAnchor}（完整四层架构）。

\noindent\textbf{任务。}
\begin{itemize}[leftmargin=1.4em]
  \item \textbf{任务1（精细对齐）：}参与者在不同头部转动速度下，使用射线或手部将虚拟标签附着到真实物体指定区域，测量完成时间与失误率。
  \item \textbf{任务2（遮挡恢复）：}交互中用书本遮挡物体，随后移开，测量参与者恢复交互信心的时间与成功率。
\end{itemize}

\noindent\textbf{主观量表。} Likert 5 级量表评估抖动、延迟、附着感、信任度四个维度（明确标注「附着/信任」为自定义改编构念），辅以 NASA-TLX~\cite{hart1988nasatlx} 认知负荷量表。

\noindent\textbf{统计。} Friedman 检验 + Wilcoxon 符号秩（Holm 校正）用于非参数主观评分；重复测量 ANOVA 用于客观指标；报告效应量（Cohen's $d$ 或 rank-biserial correlation）。伦理审查批件见 \texttt{experiment/}。

\section{实验结果}
% 写作指引：与 RQ 一一对应，先略写占位。
\subsection{空间配准与打滑消除（RQ1）}
\subsection{时序平滑与抖动抑制（RQ2）}
\subsection{断连恢复与鲁棒性（RQ3）}
\subsection{用户表现与可用性（RQ4）}

\section{讨论}

\subsection{系统贡献的本质：从感知精度到锚点质量}

本文的核心观点是：\textbf{世界一致的真实物体锚定是时间、几何、质量与恢复四个维度协同的系统问题，不能仅用单帧相机坐标系位姿精度衡量。} 计算机视觉领域的 6DoF 位姿估计以 ADD/ADD-S 等单帧度量为评价标准~\cite{hinterstoisser2012add,posecnn2017}，但这些指标刻画的是「感知有多准」，而非「锚点在运动观察者的世界坐标系中是否稳定可用」。一个在每一帧都精确的位姿流，进入混合现实渲染回路后，仍可能因时间错配而在世界空间中剧烈打滑。

EgoAnchor 的价值不在于提出新的位姿估计算法，而在于证明：\textbf{在消费级硬件上，通过四层协同的系统设计，可以将异步外部感知流转化为与平台级方案相当的锚定质量。}这为开放、跨平台的混合现实物理交互系统提供了可行路径。

\subsection{五维能力空白的系统意义}

表~\ref{tab:capability-comparison} 展示的能力矩阵不仅是技术对比，更反映了不同方案在设计空间中的取舍。带主动深度的方案（HoloLens、Vision Pro）牺牲硬件开放性换取感知质量；仅凭相机的方案（Vuforia）牺牲运动支持换取静态精度；支持运动的方案（Meta）牺牲覆盖范围换取实时性。EgoAnchor 是首个同时满足五个维度的系统，其意义在于：它证明了\textbf{在不依赖专有硬件、不限定物体类别、不牺牲运动支持的前提下，开放系统可以达到实用级别的锚定质量。}

这一结果对混合现实生态具有重要启示：未来的 MR 物理交互系统不必依赖平台封闭的 API，而可以在开放硬件上通过系统级设计达到相当的可用性。

\subsection{四层协同架构的设计原理}

EgoAnchor 的四层架构遵循「分而治之、各司其职」的设计原理：

\noindent\textbf{时间对齐层}解决物理一致性问题——帧对齐锚定确保世界空间映射对应正确的时刻，这是后续所有处理的前提。

\noindent\textbf{质量门控层}解决不确定性传播问题——拦截低质量观测，防止异常位姿污染下游模块。

\noindent\textbf{时序平滑层}解决噪声与响应折衷问题——卡尔曼状态估计提供自适应平滑，静止锁定在确定静止时完全抑制噪声。

\noindent\textbf{生命周期层}解决容错与自愈问题——追踪中断时保持有界输出，长时失效时触发重获取。

这四层的协同关键在于：每一层专注解决一个明确的子问题，层间通过清晰的接口解耦，使得整体系统既可分析、可调试，又具备端到端的协同性能。

\subsection{适用边界与失效模式}

EgoAnchor 的适用性受以下条件约束：

\noindent\textbf{物体特性。} 系统依赖视觉特征与几何一致性，对于强反光、微小无纹理或透明物体，底层感知模块会失效。当前实现假设单刚体目标，不支持可变形或铰接物体。

\noindent\textbf{运动速度。} 物体或头显过快运动（如 $>2$\,m/s 平移）导致运动模糊，视觉追踪质量下降；GPU 计算队列积压会导致时延上升，削弱帧对齐的效果。

\noindent\textbf{视野外移动。} 当物体在视野外被移动时，系统保持最后已知位姿（Coasting）或隐藏虚拟内容（Lost），但无法感知物体的新位置。这是纯视觉方案的固有限制，需要结合惯性传感器或多相机覆盖缓解。

\noindent\textbf{外部算力依赖。} 当前实现依赖外部 GPU 进行推理，网络延迟和算力波动会影响系统稳定性。未来可通过模型压缩与端侧部署降低依赖。

\subsection{局限性与未来工作}

\noindent\textbf{多物体扩展。} 当前系统聚焦单物体锚定，多物体场景需要并行追踪与独立生命周期管理，涉及算力分配与优先级调度。

\noindent\textbf{端侧轻量化。} 将感知前端部署到头显端可降低延迟并消除网络依赖，但需要针对移动 GPU 的模型压缩与加速。

\noindent\textbf{学习式时序策略。} 当前时序平滑采用手工设计的卡尔曼滤波与静止锁定，未来可探索数据驱动的时序策略学习，根据物体运动模式与用户行为自适应调整。

\noindent\textbf{长时稳定性。} 当前评估聚焦分钟级会话，长时间（小时级）运行的累积漂移、温度影响与传感器标定漂移尚未充分测试。

\section{结论}

透视混合现实中的虚拟内容与物理实体稳定锚定，需要系统同时满足多个维度的能力要求。现有方案在免主动深度、任意物体、允许运动、低成本模型与开放硬件这五个维度上各有取舍，没有任何单一系统同时覆盖。本文提出 EgoAnchor，首个填补该五维能力空白的端到端真实物体锚定系统。系统仅依赖双目透视图像与物体三维模型，运行于消费级 GPU，无需针对每个物体进行训练。

EgoAnchor 的核心是四层协同架构：时间对齐层通过帧对齐锚定消除头动打滑，质量门控层拦截异常观测，时序平滑层结合卡尔曼状态估计与静止锁定提供稳定输出，生命周期层实现追踪中断后的自愈。这四层协同设计将低频（5--12\,fps）、含噪、会间歇失效的外部感知流，转化为高频（60\,fps）、世界一致、稳定可恢复的真实物体锚点。

我们在 Quest~3 上实现并验证了该系统，建立了以锚点为中心的评估协议，从世界锚定误差、静态抖动、头动打滑、端到端时延到恢复时间全面刻画系统性能。实验表明，EgoAnchor 在消费级硬件上达到与平台级方案相当的锚定质量。系统将开源，为开放、跨平台的混合现实物理交互系统设计提供可行路径。

\nocite{*}
\bibliographystyle{abbrv-doi-hyperref-narrow}
\bibliography{egoanchor_cn_refs}

\end{document}
